% Source-derived chapter 01: Introduction
% Original source: anticyclotomic-iwasawa-rational-elliptic-eisenstein
\chapter{Introduction}
\label{anticyclotomic-iwasawa-rational-elliptic-eisenstein-chapter-01}
\source{anticyclotomic-iwasawa-rational-elliptic-eisenstein}

\begin{remark}[Source section]
\label{anticyclotomic-iwasawa-rational-elliptic-eisenstein-chapter-01-source}
\source{anticyclotomic-iwasawa-rational-elliptic-eisenstein}
\source{anticyclotomic-iwasawa-rational-elliptic-eisenstein}
This chapter reproduces the corresponding section of the retrieved
LaTeX source, including its definitions, statements, examples, and
proofs. It has not yet been translated into Lean.
\notready
\end{remark}

% The source section title was: Introduction
\addtocontents{toc}{\protect\setcounter{tocdepth}{0}}
\addtocontents{lof}{\protect\setcounter{tocdepth}{0}}
\renewcommand{\thetheorem}{\Alph{theorem}}

\subsection{Statement of the main results}

Let $E/\Q$ be an elliptic curve, and let $p$ be an odd prime of good reduction for $E$. We say that $p$ is an \emph{Eisenstein prime} (for $E$) if $E[p]$ is reducible as a $G_\Q$-module, where $G_\Q={\rm Gal}(\overline{\Q}/\Q)$ is the absolute Galois group of $\Q$ and $E[p]$ denotes the $p$-torsion of $E$. Equivalently, $p$ is an Eisenstein prime if $E$ admits a rational $p$-isogeny. By a result of Fontaine (see \cite{edixhoven-weight} for an account), Eisenstein primes are primes of \emph{ordinary} reduction for $E$, and by Mazur's results \cite{mazur} in fact $p\in\{3, 5, 7, 13, 37\}$. 
%In this paper, we study the anticyclotomic Iwasawa theory of $E$ at Eisenstein primes $p$ over an imaginary quadratic field in which $p$ splits. 

%We restrict to the elliptic curve case largely for simplicity, and it should be  noted that our methods apply more generally to the case of elliptic modular forms. 

Let $p>2$ be an Eisenstein prime for $E$, and let $K$ be an imaginary quadratic field such that 
\begin{equation}\label{eq:intro-spl}
\source{anticyclotomic-iwasawa-rational-elliptic-eisenstein}
\textrm{$p=v\bar{v}$ splits in $K$,}\tag{spl} 
\end{equation} 
where $v$ denotes the prime of $K$ above $p$ induced by a fixed embedding $\overline{\bQ}\hookrightarrow\overline{\bQ}_p$. 
%In addition, we assume that
%\begin{equation}\label{eq:intro-disc}
%\textrm{the discriminant $D_K$ of $K$ is odd and $D_K\neq -3$,}\tag{disc}
%\end{equation}
%and 
Denoting by $N$ the conductor of $E$, assume also that $K$ satisfies the following \emph{Heegner hypothesis}:
\begin{equation}\label{eq:intro-Heeg}
\source{anticyclotomic-iwasawa-rational-elliptic-eisenstein}
\textrm{every prime $\ell\vert N$ splits in $K$.}\tag{Heeg}
\end{equation}

Under these hypotheses, the anticyclotomic Iwasawa main conjecture for $E$ {considered in this paper}
can be formulated in two different guises. We begin by recalling these, since both formulations will play an important role in the proof of our main results. (Note that for the formulation $p$ can be any odd prime of good ordinary reduction for $E$.) Let $\Gamma={\rm Gal}(K_\infty/K)$ be the Galois group of the anticyclotomic $\Z_p$-extension of $K$, and for each $n$ denote by $K_n$ the subfield of $K_\infty$ with $[K_n:K]=p^n$. Set
\[
\Lambda:=\Z_p\llbracket\Gamma\rrbracket,\quad\Lambda_{\rm ac}:=\Lambda\otimes_{\Z_p}\Q_p,\quad\Lambda^{\rm ur}:=\Lambda\hat\otimes_{\Z_p}\Z_p^{\rm ur},
\]
where $\Z_p^{\rm ur}$ is the completion of the ring of integers of the maximal unramified extension of $\Q_p$. 
Following the work of Bertolini--Darmon--Prasanna \cite{BDP}, there is a $p$-adic $L$-function $\mathcal{L}_E\in\Lambda^{\rm ur}$ interpolating the central critical values of {the $L$-function of}
 $f/K$, where $f\in S_2(\Gamma_0(N))$ is the newform associated with $E$, twisted by certain characters of $\Gamma$
of infinite order. For any subfield $L\subset\overline{\Q}$, let ${\rm Sel}_{p^m}(E/L)$ be the Selmer group fitting into the descent exact sequence
\[
0\rightarrow E(L)\otimes_{\Z}\Z/p^m\Z\rightarrow{\rm Sel}_{p^m}(E/L)\rightarrow\sha(E/L)[p^m]\rightarrow 0.
\]
We put ${\rm Sel}_{p^\infty}(E/K_\infty)=\varinjlim_m{\rm Sel}_{p^m}(E/K_\infty)$, %\subset{\rm H}^1(K_\infty,E[p^\infty])$. L
and let 
\[
\mathfrak{X}_E={\rm Hom}_{\Z_p}(\mathfrak{S}_E,\Q_p/\bZ_p)
\] 
be the Pontryagin dual of {modified} Selmer group $\mathfrak{S}_E$ obtained from ${\rm Sel}_{p^\infty}(E/K_\infty)$ 
by relaxing (resp. imposing triviality) at the places above $v$ (resp. $\bar{v}$).

 
The following formulation of the anticyclotomic Iwasawa main conjectures for $E$ can be seen as a special case of Greenberg's main conjectures \cite{greenberg-motives}.

\begin{conj}
\source{anticyclotomic-iwasawa-rational-elliptic-eisenstein}
\notready\label{conj:BDP}
\source{anticyclotomic-iwasawa-rational-elliptic-eisenstein}
Let $E/\mathbb{Q}$ be an elliptic curve and $p>2$ a prime of good ordinary reduction for $E$, and let $K$ be an imaginary quadratic field satisfying {\rm (\ref{eq:intro-Heeg})} and  {\rm (\ref{eq:intro-spl})}. %, and  {\rm (\ref{eq:intro-disc})}. 
Then $\mathfrak{X}_E$ is $\Lambda$-torsion, and
\[
{\rm char}_\Lambda\bigl(\mathfrak{X}_E\bigr)\Lambda^{\rm ur}=\bigl(\Lcal_E\bigr)
\]
as ideals in $\Lambda^{\rm ur}$. 
\end{conj}

A second formulation, originally due to Perrin-Riou \cite{perrinriou}, is in terms of Heegner points. Although a more general formulation is possible (\emph{cf.} \cite[\S{2.3}]{CMpconverse}), here as in \cite{perrinriou} we assume that
\begin{equation}\label{eq:intro-disc}
\source{anticyclotomic-iwasawa-rational-elliptic-eisenstein}
\textrm{the discriminant $D_K$ of $K$ is odd and $D_K\neq -3$,}\tag{disc}
\end{equation}
and do not require hypothesis (\ref{eq:intro-spl}). Fix a modular parametrization
\[
\pi:X_0(N)\rightarrow E,
\]
and for any subfield $L\subset\overline{\Q}$ let ${\rm Sel}_{p^m}(E/L)$ be the Selmer group fitting into the descent exact sequence
\[
0\rightarrow E(L)\otimes_{\Z}\Z/p^m\Z\rightarrow{\rm Sel}_{p^m}(E/L)\rightarrow\sha(E/L)[p^m]\rightarrow 0.
\]
%where $\sha(E/L)$ is the Tate--Shafarevich group of $E/L$. 
Via $\pi$, the Kummer images of Heegner points on $X_0(N)$ over ring class fields of $K$ of $p$-power conductor give rise to a class $\kappa_1^{\rm Hg}\in\mathcal{S}$, where
\[
\mathcal{S}=\biggl(\varprojlim_n\varprojlim_m{\rm Sel}_{p^m}(E/K_n)\biggr)\otimes_{}\Q_p.
\] 
{The group $\mathcal{S}$ is naturally a $\Lambda_{\rm ac}$-module, and} the class $\kappa_1^{\rm Hg}$ is known to be non-$\Lambda_{\rm ac}$-torsion by results of Cornut and Vatsal \cite{cornut}, \cite{vatsal}. Denote by $\mathcal{H}\subset\mathcal{S}$ the $\Lambda_{\rm ac}$-submodule generated by $\kappa_1^{\rm Hg}$, and put 
\[
\mathcal{X}={\rm Hom}_{\Z_p}({\rm Sel}_{p^\infty}(E/K_\infty),\Q_p/\Z_p)\otimes\Q_p.
\]
%The first formulation of the anticyclotomic main conjecture for $E$ was given by Perrin-Riou \cite{perrinriou}:

%Let $f\in S_2(\Gamma_0(N))$ be the newform associated to $E$, so that $L(f,s)=L(E,s)$, and let $c_E\in\Z$ be such that $\pi^*(\omega_E)=c_E\cdot 2\pi i f(\tau)d\tau$, for $\omega_E$ a N\'{e}ron differential on $E$.

\begin{conj}
\source{anticyclotomic-iwasawa-rational-elliptic-eisenstein}
\notready\label{conj:PR}
\source{anticyclotomic-iwasawa-rational-elliptic-eisenstein}
Let $E/\Q$ be an elliptic curve and $p>2$ a prime of good ordinary reduction for $E$, and let $K$ be an imaginary quadratic field satisfying {\rm(\ref{eq:intro-Heeg})} and {\rm(\ref{eq:intro-disc})}. Then $\mathcal{S}$ and $\mathcal{X}$ both
have $\Lambda_{\rm ac}$-rank one, and
\[
{\rm char}_{\Lambda_{\rm ac}}(\mathcal{X}_{\rm tors})={\rm char}_{\Lambda_{\rm ac}}\bigl(\mathcal{S}/\mathcal{H}\bigr)^2,
\] 
where $\mathcal{X}_{\rm tors}$ denotes the $\Lambda_{\rm ac}$-torsion submodule of $X$. 
\end{conj}

\begin{rem}
\source{anticyclotomic-iwasawa-rational-elliptic-eisenstein}
\notready
One can naturally formulate an integral version of Conjecture~\ref{conj:PR}, but the results of \cite{perrinriou} and \cite{howard-GZ} show that the terms appearing in the corresponding equality of $\Lambda$-module characteristic ideals are in general \emph{not} invariant under isogenies. (With $p$ inverted, i.e., as ideals in $\Lambda_{\rm ac}$, the terms are invariant under isogenies.) On the other hand, it is clear that the principal ideals in $\Lambda^{\rm ur}$ appearing in the equality of Conjecture~\ref{conj:BDP} depend only on the isogeny class of $E$.
\end{rem}

When $p$ is non-Eisenstein for $E$, Conjectures~\ref{conj:BDP} and \ref{conj:PR} have been studied by several authors \cite{bertolini,howard,howard-gl2-type,wan-heegner,castellaheights,BCK}, but the residually reducible case remained largely unexplored; in particular, unless $E$ has CM by $K$ (a case that is excluded by our hypothesis (\ref{eq:intro-Heeg}), but see \cite{CMpconverse} for this case), there seems to be no previous results towards these conjectures when $p$ is an Eisenstein prime for $E$. 

To state our main results on %Conjectures~\ref{conj:PR} and \ref{conj:BDP} 
the anticyclotomic Iwasawa theory of $E$ at Eisenstein primes $p$, write
\[
E[p]^{ss}=\mathbb{F}_p(\phi)\oplus\mathbb{F}_p(\psi),
\]
where $\phi,\psi:G_\Q\rightarrow\mathbb{F}_p^\times$ are characters. Note that it follows from the Weil pairing that $\psi=\omega\phi^{-1}$, where $\omega$ is the Teichm\"uller character. Let $G_p\subset G_{\mathbb{Q}}$ be a decomposition group at $p$. 

Our most complete results towards Conjectures \ref{conj:BDP} and \ref{conj:PR} are proved under the additional hypothesis
that
\begin{equation}\label{eq:intro-sel}
\source{anticyclotomic-iwasawa-rational-elliptic-eisenstein}
\text{the $\Z_p$-corank of $\mathrm{Sel}_{p^\infty}(E/K)$ is $1$.}\tag{Sel}
\end{equation}
%\textcolor{red}{Shall we add here the hypothesis
%\begin{equation}\label{eq:intro-CM}
%E \text{ does not have CM by } K. \tag{K-CM}
%\end{equation}}

\begin{theorem}
\source{anticyclotomic-iwasawa-rational-elliptic-eisenstein}
\notready\label{thm:C}
\source{anticyclotomic-iwasawa-rational-elliptic-eisenstein}
Let $E/\mathbb{Q}$ be an elliptic curve, $p>2$ an Eisenstein prime for $E$, and $K$ an imaginary quadratic field satisfying {\rm (\ref{eq:intro-Heeg})}, {\rm (\ref{eq:intro-spl})}, {\rm (\ref{eq:intro-disc})}, and
{\rm (\ref{eq:intro-sel})}. 
Assume also that
$\phi\vert_{G_p}\neq\mathds{1},\omega$. Then Conjecture~\ref{conj:BDP} holds.
\end{theorem}


As we explain in more detail in the next section, a key step towards Theorem~\ref{thm:C} is the proof of a divisibility in  Conjecture~\ref{conj:PR} that we establish without the need to assume (\ref{eq:intro-spl}) (see Theorem~\ref{thm:howard-HP}). On the other hand, as first observed in \cite{castella-PhD} and \cite{wan-heegner}, when $p$ splits in $K$, Conjectures~\ref{conj:BDP} and \ref{conj:PR} are essentially equivalent (see Proposition \ref{prop:equiv-imc}). Thus from Theorem~\ref{thm:howard-HP} we deduce one of the divisibilities in Conjecture~\ref{conj:BDP}, which by the analysis of Iwasawa invariants carried out in $\S\S$\ref{sec:algebraic}-\ref{IMCII} then yields the equality of ideals in $\Lambda^{\rm ur}$ predicted by Conjecture~\ref{conj:BDP}. As a result, our analysis together with the aforementioned equivalence also yields the following.


\begin{cor}
\source{anticyclotomic-iwasawa-rational-elliptic-eisenstein}
\notready\label{cor:D}
\source{anticyclotomic-iwasawa-rational-elliptic-eisenstein}
%Under the hypotheses of Theorem~\ref{thm:C}, both $\mathcal{S}$ and $X$ have $\Lambda_{\rm ac}$-rank one, and
%\[
%{\rm char}_{\Lambda_{\rm ac}}(X_{\rm tors})={\rm char}_{\Lambda_{\rm ac}}\bigl(\mathcal{S}/\mathcal{H}\bigr)^2.
%\]
Let $E/\mathbb{Q}$ be an elliptic curve, $p>2$ an Eisenstein prime for $E$, and $K$ an imaginary quadratic field satisfying {\rm (\ref{eq:intro-Heeg})}, {\rm (\ref{eq:intro-disc})}, and
{\rm (\ref{eq:intro-sel})}. If $E(K)[p]=0$, then $\mathcal{S}$ and $\mathcal{X}$ both
have $\Lambda_{\rm ac}$-rank one, and
\[
{\rm char}_{\Lambda_{\rm ac}}(\mathcal{X}_{\rm tors})\supset{\rm char}_{\Lambda_{\rm ac}}\bigl(\mathcal{S}/\mathcal{H}\bigr)^2.
\] 
Moreover, if in addition $K$ satisfies {\rm (\ref{eq:intro-spl})} and $\phi\vert_{G_p}\neq\mathds{1},\omega$, then
\[
{\rm char}_{\Lambda_{\rm ac}}(\mathcal{X}_{\rm tors})={\rm char}_{\Lambda_{\rm ac}}\bigl(\mathcal{S}/\mathcal{H}\bigr)^2,
\] 
and hence Conjecture~\ref{conj:PR} holds.
\end{cor}

%\textcolor{red}{More precisely, as explained below, we prove the vanishing of $\mu$-invariants and equality of $\lambda$-invariants in Conjecture~B, which reduces Theorem \ref{thm:C} to proving one divisibility of characteristic ideals. Our strategy exploits the relation between these two conjectures (see Proposition \ref{prop:equiv-imc}). We indeed prove one divisibility in Conjecture~A using a Kolyvagin system argument and then obtain both Theorem \ref{thm:C} and Corollary \ref{cor:D}.}

Note that in both Theorem~\ref{thm:C} and Corollary~\ref{cor:D}, the elliptic curve $E$ is allowed to have complex multiplication  (necessarily by an imaginary quadratic field different from $K$). 

With a judicious choice of $K$, Theorem~\ref{thm:C} also has applications to the arithmetic over $\Q$ of rational elliptic curves. Specifically, for Eisenstein prime $p$, we obtain a $p$-converse to the celebrated theorem 
\begin{equation}\label{eq:GZK}
\source{anticyclotomic-iwasawa-rational-elliptic-eisenstein}
\ord_{s=1}L(E,s)=r \in \{0,1\}\;\Longrightarrow\;
{\rm rank}_\Z E(\Q)=r\;\textrm{and}\;
\#\sha(E/\bQ)<\infty,
%\#\sha(E/\bQ)[p^\infty]<\infty,
\end{equation}
of Gross--Zagier and Kolyvagin. (The case of Eisenstein primes eluded the methods of 
%\cite{pCONVskinner,wei-zhang-mazur-tate,jsw}
\cite{pCONVskinner} and \cite{wei-zhang}, which require $E[p]$ to be absolutely irreducible as a $G_\bQ$-module.)

\begin{theorem}
\source{anticyclotomic-iwasawa-rational-elliptic-eisenstein}
\notready\label{thm:E}
\source{anticyclotomic-iwasawa-rational-elliptic-eisenstein}
Let $E/\mathbb{Q}$ be an elliptic curve and $p>2$ an Eisenstein prime for $E$, so that $E[p]^{ss}=\mathbb{F}_p(\phi)\oplus\mathbb{F}_p(\psi)$ as $G_\Q$-modules, and assume that $\phi\vert_{G_p}\neq\mathds{1},\omega$. Let $r\in\{0,1\}$. Then the following implication holds:
\[
%\left.
%\begin{array}{r}
%{\rm rank}_\Z E(\bQ)=r\\[0.2em]
%\#\sha(E/\bQ)[p^\infty]<\infty
%\end{array}
%\right\}
{\rm corank}_{\Z_p}{\rm Sel}_{p^\infty}(E/\Q)=r
\;\Longrightarrow\;
{\rm ord}_{s=1}L(E,s)=r,
\]
and so ${\rm rank}_\Z E(\Q)=r$ and $\#\sha(E/\Q)<\infty$.
\end{theorem}

Note that if ${\rm rank}_{\Z}E(\Q)=r$ and $\#\sha(E/\Q)[p^\infty]<\infty$ then ${\rm corank}_{\Z_p}{\rm Sel}_{p^\infty}(E/\Q)=r$,
whence the $p$-converse to \eqref{eq:GZK}. We also note that for $p=3$, Theorem~\ref{thm:E} together with the recent work of Bhargava--Klagsbrun--Lemke Oliver--Shnidman \cite{BKLS} on the average $3$-Selmer rank for abelian varieties in quadratic twist families, 
%improves on the best known results
provides additional evidence
towards Goldfeld's conjecture \cite{goldfeld} for elliptic curves $E/\bQ$ admitting a rational $3$-isogeny (see Corollary~\ref{cor:goldfeld} and Remark~\ref{rem:BKLS}, and see also \cite{kriz-li} for earlier results along these lines).

Another application of Theorem~\ref{thm:C} is the following.
\begin{theorem}
\source{anticyclotomic-iwasawa-rational-elliptic-eisenstein}
\notready\label{thm:F}
\source{anticyclotomic-iwasawa-rational-elliptic-eisenstein}
Under the hypotheses of Theorem~\ref{thm:E}, assume in addition that $\phi$ is either ramified at $p$ and odd, or unramified at $p$ and even. If ${\rm ord}_{s=1}L(E,s)=1$, then 
\[
\ord_{p}\biggl(\frac{L'(E, 1)}{\operatorname{Reg}(E / \Q) \cdot \Omega_{E}}\biggr)=\ord_{p}\biggl(\# \sha(E / \Q) \prod_{\ell \nmid \infty} c_{\ell}(E /\Q)\biggr),
\]
where 
\begin{itemize}
\item ${\rm Reg}(E/\Q)$ is the regulator of $E(\Q)$, 
\item $\Omega_E=\int_{E(\mathbb{R})}\vert\omega_E\vert$ is the N\'{e}ron  period associated to the N\'{e}ron differential $\omega_E$, and
\item $c_\ell(E/\Q)$ is the Tamagawa number of $E$ at the prime $\ell$.
\end{itemize} 
In other words, the $p$-part of the Birch--Swinnerton-Dyler formula for $E$ holds.
\end{theorem}


\subsection{Method of proof and outline of the paper}

Let us explain some of the ideas that go into the  proof of our main results, beginning with the proof of Theorem~\ref{thm:C}. As in \cite{greenvats}, our starting point is Greenberg's old observation \cite{greenberg-nagoya} that a ``main conjecture'' should be equivalent to an imprimitive one. More precisely, in the context of Theorem~\ref{thm:C}, for $\Sigma$ any finite set of non-archimedean primes of $K$ not containing any of the primes above $p$, this translates into the expectation that the $\Sigma$-imprimitive Selmer group $\mathfrak{X}_E^\Sigma$, obtained by relaxing the local condition defining (the Pontryagin dual of) $\mathfrak{X}_E$ at the primes $w\in\Sigma$, is $\Lambda$-torsion with
\begin{equation}\label{eq:BDP-imp}
\source{anticyclotomic-iwasawa-rational-elliptic-eisenstein}
{\rm char}_\Lambda\bigl(\mathfrak{X}_E^\Sigma\bigr)\Lambda^{\rm ur}\overset{?}=\bigl(\mathcal{L}_E^\Sigma\bigr)
\end{equation}
as ideals in $\Lambda^{\rm ur}$, where $\mathcal{L}_E^\Sigma:=\mathcal{L}_E\cdot\prod_{w\in\Sigma}\mathcal{P}_w(E)$ for certain elements in $\mathcal{P}_w(E)\in\Lambda$ interpolating, for varying characters $\chi$ of $\Gamma$, the $w$-local Euler factor of $L(E/K,\chi,s)$ evaluated $s=1$. 

A key advantage of the imprimitive main conjecture (\ref{eq:BDP-imp}) is that (unlike the original conjecture), for suitable choices of $\Sigma$, its associated Iwasawa invariants are well-behaved with respect to congruences mod $p$. Identifying $\Lambda$ with the power series ring $\Z_p\llbracket T\rrbracket$ by setting $T=\gamma-1$ for a fixed topological generator $\gamma\in\Gamma$, recall that by the Weierstrass preparation theorem, 
every nonzero $g\in\Lambda$ can be uniquely written in the form
\[
g=u\cdot p^\mu\cdot Q(T),
\]
with $u\in\Lambda^\times$, $\mu=\mu(g)\in\Z_{\geqslant 0}$, and $Q(T)\in\Z_p[T]$ a distinguished polynomial of degree $\lambda(g)$. The constants $\lambda$ and $\mu$ are the so-called \emph{Iwasawa invariants} of $g$. For a torsion $\Lambda$-module $\mathfrak{X}$ we let $\lambda(\mathfrak{X})$ and $\mu(\mathfrak{X})$ be the Iwasawa invariants of a characteristic power series for $\mathfrak{X}$, and for a nonzero $\mathcal{L}\in\Lambda^{\rm ur}$ we let $\lambda(\Lcal)$ and $\mu(\Lcal)$ be the Iwasawa invariants of any element of $\Lambda$ generating the same $\Lambda^{\rm ur}$-ideal as $\mathcal{L}$. 

As a first step towards Theorem~\ref{thm:C}, we deduce from the $G_\Q$-module isomorphism $E[p]^{ss}=\mathbb{F}_p(\phi)\oplus\mathbb{F}_p(\psi)$ that, taking $\Sigma$ to consist of primes that are split in $K$ and containing all the primes of bad reduction for $E$, the module $\mathfrak{X}_E^\Sigma$ is $\Lambda$-torsion with
\begin{equation}\label{eq:alg-inv}
\source{anticyclotomic-iwasawa-rational-elliptic-eisenstein}
\mu\bigl(\mathfrak{X}_E^\Sigma\bigr)=0\quad\textrm{and}\quad\lambda\bigl(\mathfrak{X}_E^\Sigma\bigr)=\lambda\bigl(\mathfrak{X}_\phi^\Sigma\bigr)+\lambda\bigl(\mathfrak{X}_\psi^\Sigma\bigr),
\end{equation}
where $\mathfrak{X}_\phi^\Sigma$ and $\mathfrak{X}_\psi^\Sigma$ are anticyclotomic Selmer groups (closely related to the Pontryagin dual of certain class groups) for the Teichm\"uller lifts of $\phi$ and $\psi$, respectively. The proof of (\ref{eq:alg-inv}), which is taken up in $\S\ref{sec:algebraic}$, uses Rubin's work \cite{rubinmainconj} on the Iwasawa main conjecture for imaginary quadratic fields and Hida's work \cite{hidamu=0} on the vanishing of the $\mu$-invariant of $p$-adic Hecke $L$-functions. 

On the other hand, in $\S\ref{IMCII}$ we deduce from the main result of \cite{kriz} that for such $\Sigma$ one also has
\begin{equation}\label{eq:an-inv}
\source{anticyclotomic-iwasawa-rational-elliptic-eisenstein}
\mu\bigl(\Lcal_E^\Sigma\bigr)=0\quad\textrm{and}\quad\lambda\bigl(\Lcal_E^\Sigma\bigr)=\lambda\bigl(\Lcal_\phi^\Sigma\bigr)+\lambda\bigl(\Lcal_\psi^\Sigma\bigr),
\end{equation}
where $\mathcal{L}_\phi^\Sigma$ and $\Lcal_\psi^\Sigma$ are $\Sigma$-imprimitive anticyclotomic Katz $p$-adic $L$-functions attached to $\phi$ and $\psi$, respectively. 

With equalities (\ref{eq:alg-inv}) and (\ref{eq:an-inv}) in hand, 
it follows easily that to prove  the equality of characteristic ideals in Conjecture~\ref{conj:BDP} it suffices to prove one of the predicted divisibilities in $\Lambda^{\rm ur}[\frac{1}{p}]$. In $\S\ref{IMCIII}$,
{by combining Howard's approach to proving Iwasawa-theoretic divisibilities \cite{howard} with a Kolyvagin system argument along the lines of 
Nekov\'a\v{r}'s \cite{nekovar} (but adapted for twists by infinite order characters and for obtaining a bound on the length of Tate--Shafarevich group and not just an annihilator),}
% by an extension of Howard's Kolyvagin system arguments \cite{howard} to the residually reducible setting, 
we prove the main result towards one of the divisibilities in Conjecture~\ref{conj:PR}: 
{${\rm char}_{\Lambda_{\rm ac}}(\mathcal{X}_{\rm tors})$ divides ${\rm char}_{\Lambda_{\rm ac}}\bigl(\mathcal{S}/\mathcal{H}\bigr)^2$ in 
$\Lambda_{\rm ac}[\frac{1}{T}]$, and even in 
$\Lambda_{\rm ac}$ assuming \eqref{eq:intro-sel}.} As already noted, hypotheses (\ref{eq:intro-spl}) and $\phi\vert_{G_p}\neq\mathds{1},\omega$ are not needed at this point. 
This yields a corresponding divisibility in (\ref{eq:BDP-imp}), from which the proof of Theorem~\ref{thm:C} follows easily. The details of the final argument, and the deduction of Corollary~\ref{cor:D}, are given in $\S\ref{sec:CD}$.
The additional hypothesis \eqref{eq:intro-sel} is required to circumvent the growth of the
`error term' in our 
%extension of Howard's 
Kolyvagin system arguments in the cases of twists by anticyclotomic characters 
$p$-adically close to the trivial character. {The arguments in $\S\ref{IMCIII}$ apply equally well to both the residually reducible and residually irreducible cases.}

Finally, the proofs of Theorems~\ref{thm:E} and \ref{thm:F} are given in $\S\ref{sec:EF}$, %and similarly as in \cite{pCONVskinner} and \cite{jsw} 
and they are both obtained as an application of Theorem~\ref{thm:C} for a suitably chosen $K$. In particular, the proof of Theorem~\ref{thm:F} requires knowing the $p$-part of the Birch--Swinnerton-Dyler formula in analytic rank $0$
for the quadratic twist $E^K$; this is deduced in Theorem~\ref{thmGV} from the results of Greenberg--Vatsal \cite{greenvats}, and this is responsible for the additional hypotheses on $\phi$ placed in Theorem~\ref{thm:F}. 
%The removal of these hypotheses is the subject of subsequent work.


\subsection{Examples}

To illustrate Theorem~\ref{thm:F}, take $p=5$ and consider the elliptic curve
\begin{displaymath}
J: y^2+y=x^3+x^2-10x+10.
\end{displaymath} 
The curve $J$ has conductor $123$ and analytic rank $1$, and satisfies $J[5]^{ss}=\Z/5\Z\oplus\mu_p$ as $G_\Q$-modules ($J$ has a rational $5$-torsion point). If $\psi$ is an even quadratic character such that $\psi(5)=-1$, corresponding to a real quadratic field $\Q(\sqrt{c})$ in which $5$ is inert, then the twist $E=J_c$ of $J$ by $\psi$ satisfies the hypotheses of Theorem~\ref{thm:E} with $p=5$. Since the root number of $J$ is $-1$ (being of analytic rank one), by \cite[Thm.~B.2]{twist} we can find infinitely many $\psi$ as above for which the associated twist $E=J_c$ also has analytic rank one, and therefore for which Theorem~\ref{thm:F} applies. 

One can proceed similarly for $p=3$ (resp. $p=7$), taking real quadratic twists of, for example, the elliptic curve $y^2+y=x^3+x^2-7x+5$ of conductor 91 (resp. $y^2+xy+y=x^{3}-x^{2}-19353x+958713$ of conductor 574). For $p=13$ (resp. $p=37$), one can do the same, possibly choosing the quadratic character to be odd and/or imposing conditions at some bad primes depending on the character describing the kernel of the isogeny (which is not trivial in these cases) in order to apply \cite[Thm.~B.2]{twist}. One could consider, for example, twists of the elliptic curve $y^2+y=x^3-21x+40$ of conductor 441 (resp. $y^2+xy+y=x^3+x^2-8x+6$ of conductor 1225). 

We also note that, for each of the four primes primes above, $p=3,5,7,13$, there are infinitely many distinct $j$-invariants to which Theorem~\ref{thm:F} applies, as $X_0(p)$ has genus $0$ in these cases. 


\subsection{Relation to previous works}

Results in the same vein as (\ref{eq:alg-inv}) and (\ref{eq:an-inv}) were first obtained by Greenberg--Vatsal \cite{greenvats} in the cyclotomic setting; combined with Kato's Euler system divisibility \cite{kato-euler-systems}, these results led to their proof of the cyclotomic Iwasawa main conjecture for rational elliptic curves at Eisenstein primes $p$ (under some hypotheses on the kernel of the associated rational $p$-isogeny). This paper might be seen as an extension of the Greenberg--Vatsal method for Eisenstein primes to the anticyclotomic setting.  
{However, for the anticyclotomic Selmer groups and $L$-functions considered in this paper we are able to avoid the possible variation within an isogeny class of elliptic curves of the $\mu$-invariants and periods, which must be dealt with in \cite{greenvats}. In large part this is because the periods
in the corresponding $p$-adic families are the CM periods of Hecke characters and not the periods of the elliptic curve.  Consequently, the methods are slightly more robust and the resulting applications somewhat more general. 
The $\mu$-invariants of anticyclotomic Selmer groups and $p$-adic $L$-functions were also studied in \cite{pollack2011anticyclotomic}, but for different Selmer conditions and hypotheses on $K$ (in fact, under the hypothesis \eqref{eq:intro-Heeg} the $p$-adic $L$-function in \cite{pollack2011anticyclotomic} vanishes identically and the Selmer group is not $\Lambda$-cotorsion).}


The ensuing applications {of Theorems \ref{thm:C} and Corollary \ref{cor:D}} to the $p$-converse of the Gross--Zagier--Kolyvagin theorem (Theorem~\ref{thm:E}) and the $p$-part of the Birch--Swinnerton-Dyer formula in analytic rank $1$ (Theorem~\ref{thm:F}) covers primes $p$ that were either left untouched by the recent works in these directions \cite{pCONVskinner,pCONVvenerucci,pBSDbert,wei-zhang,skinner-zhang,jsw,castella,pBSDbps} (where $p$ is assumed to be non-Eisenstein), or extending previous works \cite{tian,CLTZ,CCL} ($p=2$),  \cite{kriz-li} ($p=3$), \cite{CMpconverse} (CM cases). {Many of these results (especially \cite{pCONVskinner} and \cite{jsw}) also rely on progress toward 
 Conjecture \ref{conj:BDP} in the residually irreducible case. Such progress has generally come via Eisenstein congruences on higher rank unitary groups and has explicitly excluded the Eisenstein cases considered in this paper.}


\subsection{Weight two newforms} The methods and results of this paper should easily extend to cuspidal newforms of weight two and trivial character that are congruent to Eisenstein series at a prime above $p$. We have focused on the case of elliptic curves in the interest of not obscuring the main features of our argument with cumbersome notation. The general case will be addressed in later work that will also consider higher weight forms as well as Hilbert modular forms.

\subsection{Acknowledgements}
This paper has its origins in one of the projects proposed by F.C. and C.S. at the 2018 Arizona Winter School on Iwasawa theory, and we would like to thank the organizers for making possible a uniquely stimulating week during which G.G. and J.L. made very substantial progress on this project. 
{We also thank Ashay Burungale and the anonymous referees for many helpful comments and suggestions on an earlier draft of this paper.}
G.G. is grateful to Princeton University for the hospitality during a visit to F.C. and C.S. in February-March 2019. During the preparation of this paper, F.C. was partially supported by the NSF grant DMS-1946136 and DMS-2101458; G.G. was supported by the Engineering and Physical Sciences Research Council [EP/L015234/1], the EPSRC Centre for Doctoral Training in Geometry and Number Theory (The London School of Geometry and Number Theory), University College London;
C.S. was partially supported by the Simons Investigator Grant \#376203 from the Simons Foundation and by the NSF grant DMS-1901985.


\addtocontents{toc}{\protect\setcounter{tocdepth}{2}}
\addtocontents{lof}{\protect\setcounter{tocdepth}{2}}

\renewcommand{\thetheorem}{\arabic{section}.\arabic{subsection}.\arabic{theorem}}
